An \openshmem nonblocking collective operation, like a blocking collective
operation, is a group communication operation among the
participants of the team. All \acp{PE} in the team are required to call the
collective operation and each collective operation must be initiated in the same
program order across all \acp{PE}. The execution of the collective operations
may be performed in any order when the collective operation updates buffers that
are not in use as defined in Section \ref{subsec:invoking_openshmem_operations}.
Collective operations that perform updates to in use buffers will be executed in order.

\begin{enumerate}

\item Invocation semantics: Upon invocation of a nonblocking collective routine,
the operation is initiated and the routine returns without ensuring completion. All \acp{PE} in the team
must call this routine with identical arguments.

\item Collective Types: The nonblocking variants supported include the alltoall
and broadcast collectives. All other collective operations such as
reductions, collect, fcollect, barrier, barrier all, alltoalls, sync, and sync all will not have nonblocking variants.

\item Completion semantics: \openshmem programs can learn the status of the collective operations
using the \FUNC{shmem\_req\_test} routine. Completion of the operation can be
observed through one or more calls to \FUNC{shmem\_req\_test} or a single
call to \FUNC{shmem\_req\_wait}.

\item Threads: While using SHMEM\_THREAD\_MULTIPLE, the \openshmem
programs are not allowed to call multiple collective operations on different threads
and the same team.

\end{enumerate}

Note: Like other nonblocking \openshmem operations, the implementations are
expected to asynchronously progress the collective operations. The guidance on
asynchronous progress is provided in Section \ref{subsec:progress}. Additionally,
the use of \FUNC{shmem\_quiet} and \FUNC{shmem\_barrier\_all} do not guarantee
completion of a nonblocking collective operation.
